%%% Estructura de la tesis %%%
\pdfbookmark[0]{Structure}{Structure}
\chapter*{Estructura de la Tesis Doctoral}

Esta Tesis Doctoral está estructurada en seis secciones. El \textbf{Capítulo \ref{chap:introduction}} presenta una introducción del trabajo, dividida en dos secciones: Microbioma y Métodos y Técnicas Computacionales. La primera sección establece el marco biológico del estudio, mientras que la segunda sección introduce las herramientas y enfoques que se han empleado a lo largo de esta Tesis Doctoral. En el \textbf{Capítulo \ref{chap:objectives}} se enumeran los objetivos propuestos, los cuales han motivado la realización de esta Tesis Doctoral.

El \textbf{Capítulo \ref{chap:results}} expone los resultados obtenidos durante el desarrollo de la Tesis Doctoral. Este capítulo está dividido en tres secciones, las cuales se corresponden con los trabajos que conforman el compendio. En cada una de ellas se expone el contexto y los resultados obtenidos en cada uno de los trabajos. Además, se adjunta una copia de cada uno de ellos. Las referencias de los trabajos, por orden de aparición, son las siguientes:

\begin{my_list_item}
	
    \item Loaiza, L. I. S., \textbf{Fernández-Edreira, D.}, Liñares-Blanco, J., Cepeda, A., Cardelle-Cobas, A., and Fernandez-Lozano, C. (2025). Fecal microbiome analysis in patients with metabolic syndrome and type 2 diabetes. PeerJ, 13, e19108. Q2, 2.4 IF. 1 cita (Google Scholar). \url{https://doi.org/10.7717/peerj.19108}
    
    \item \textbf{Fernández-Edreira, D.}, Liñares-Blanco, J., and Fernandez-Lozano, C. (2021). Machine Learning analysis of the human infant gut microbiome identifies influential species in type 1 diabetes. Expert Systems with Applications, 185, 115648. Q1, 8.665 IF. 48 citas (Google Scholar). \url{https://doi.org/10.1016/j.eswa.2021.115648}
    
    \item \textbf{Fernández-Edreira, D.}, Liñares-Blanco, J., Patricia, V., and Fernandez-Lozano, C. (2024). VIBES: A consensus subtyping of the vaginal microbiota reveals novel classification criteria. Computational and Structural Biotechnology Journal, 23, 148-156. Q2, 4.1 IF. 4 citas (Google Scholar). \url{https://doi.org/10.1016/j.csbj.2023.11.050}
    
\end{my_list_item}

En el \textbf{Capítulo \ref{chap:discussion}}, se lleva a cabo una discusión integrada de los tres trabajos realizados, resaltando los puntos comunes que fueron esenciales en su desarrollo y brindando coherencia, cohesión y justificación al trabajo de investigación en su conjunto. El \textbf{Capítulo \ref{chap:conclussions}} enumera las conclusiones obtenidas a lo largo del desarrollo de estos tres trabajos. Finalmente, en el \textbf{Capítulo \ref{chap:future}} se exponen los futuros desarrollos que han surgido tras la investigación realizada a lo largo de esta Tesis Dcotoral.

